\needspace{10\baselineskip}
\begin{center}
\pdfbookmark[1]{9. Zu Art. 138.}{zusatz-009}
[9. Zu Art. 138.]
\end{center}
Das Criterium des Falls, wo der die geocentrische Bewegung tangirende grösste Kreis durch den Sonnenort geht, lässt sich sehr elegant auf die heliocentrischen Positionen beziehen.

Es seien
\begin{center}
\begin{tabular}{@{}l@{\quad}p{.79\linewidth}@{}}
\(R,r\)&Abstände von der Sonne für Erde und Planet;\\
\(L,l\)&Längen der Erde und des Planeten, letztere in der Bahn und beide vom aufsteigenden Knoten gezählt;\\
\(E,e\)&Excentricitäten;\\
\(V,v\)&wahre Anomalien;\\
\(P,p\)&halbe Parameter;
\end{tabular}
\end{center}
dann ist das Criterium jenes Falls in der Gleichung
\[
\frac{r\sin l}{\BareRoot{p}}=\frac{R\sin L}{\BareRoot{P}}
\quad\text{oder}\quad
\frac{\BareRoot{p}\mathbin{.}\sin l}{1+e\cos v}=\frac{\BareRoot{P}\mathbin{.}\sin L}{1+E\cos V}
\]
enthalten.

\SourceLeaf{310}
Der Beweis stützt sich auf folgende Figur:
\begin{center}
\begin{tikzpicture}[x=.35pt,y=-.35pt,line width=.6pt,line cap=round,line join=round]
\path[use as bounding box] (1640,600) rectangle (2440,1010);
\draw (1665,818)..controls(1720,816)and(1790,816)..(1855,817)..controls(2040,815)and(2260,840)..(2371,865);
\draw (1678,952)..controls(1755,951)and(1820,946)..(1887,937)..controls(2070,927)and(2260,886)..(2371,865);
\draw (1796,632)--(1855,817)--(1887,937);
\fill (1796,632) circle[radius=2.5pt];
\fill (1800,816.5) circle[radius=2.5pt];
\fill (1835,944) circle[radius=2.5pt];
\node[inner sep=0pt,font=\fontsize{10}{11}\selectfont] at (1822,632) {\(\mathsf{B}\)};
\node[inner sep=0pt,font=\fontsize{10}{11}\selectfont] at (1805,781) {\(\mathsf{C}'\)};
\node[inner sep=0pt,font=\fontsize{10}{11}\selectfont] at (1870,782) {\(\mathsf{C}\)};
\node[inner sep=0pt,font=\fontsize{10}{11}\selectfont] at (2081,785) {\(\mathsf{l}\)};
\node[inner sep=0pt,font=\fontsize{10}{11}\selectfont] at (1833,982) {\(\mathsf{A}'\)};
\node[inner sep=0pt,font=\fontsize{10}{11}\selectfont] at (1890,980) {\(\mathsf{A}\)};
\node[inner sep=0pt,font=\fontsize{10}{11}\selectfont] at (2075,951) {\(\mathsf{L}\)};
\node[inner sep=0pt,font=\fontsize{10}{11}\selectfont] at (2400,872) {\(\AscendingNode\)};
\end{tikzpicture}
\end{center}

\(A\) Ort der Erde, \(B\) geoc., \(C\) helioc. Ort des Planeten; \(A',C'\) Plätze von \(A\) und \(C\) nach unendl. kl. Zwischenzeit \(\mathrm{d}t\).

Es ist also
\[
AA'=\mathrm{d}L,\qquad CC'=\mathrm{d}l.
\]
Vermöge der Natur des Problems müssen die beiden neuen Plätze von der durch Sonne, Erde und Planet gelegten (durch \(ACB\) repräs. Ebene) gleich weit abstehen. Diese Abstände sind \(R\mathrm{d}L\sin A\) und \(r\mathrm{d}l\sin C\); da nun \(\sin A:\sin C=\sin l:\sin L\), so wird \(\dfrac{R\mathrm{d}L}{\sin L}=\dfrac{r\mathrm{d}l}{\sin l}\mathpunct{\raisebox{.45ex}{,}}\) oder da \(RR\mathrm{d}L:rr\mathrm{d}l=\BareRoot{P}:\BareRoot{p}\):
\[
\frac{\BareRoot{P}}{R\sin L}=\frac{\BareRoot{p}}{r\sin l}\mathpunct{\raisebox{.45ex}{,}}
\]
welches die obige Gleichung ist.
\begin{center}\rule{.18\textwidth}{.35pt}\end{center}
